\section{FrameWorkers-Bench: Sub-agent Routing Benchmark}
\label{sec:bench}

We evaluate director routing as a standalone task: given a natural-
language user goal $g$ for a video deliverable, produce the ordered
chain of sub-agents whose composition satisfies it. Routing is
combinatorial over a closed vocabulary of $20$ agents and multi-step
(plans range from $2$ to $10$ agents); routing errors compound
downstream because each agent consumes specific upstream artifacts.

\paragraph{Agents.}
The inventory $\mathcal{A}$ contains $20$ agents across modalities:
text (Story, Screenplay, Narration, Translation), image
(Illustration, KeyFrame, IntakeImage, BriefEnricher), video (Video,
IntakeVideo, VideoAnalysis, StyleTransfer, VideoExtend, Highlight),
audio (Music, Ambience, Narrator, AudioMix, Transcription), and
final composition (Compositor). Each agent carries a typed descriptor
specifying its \texttt{Inputs}, \texttt{Output}, and
\texttt{Purpose / Trigger}; the descriptors are concatenated verbatim
into the planner's system prompt.

\paragraph{Task.}
A planner $\pi$ maps $g$ to an ordered chain
$\hat{\mathbf{c}} = (\hat{c}_1, \dots, \hat{c}_n)$ with
$\hat{c}_i\!\in\!\mathcal{A}$. Hard structural constraints: each
$\hat{c}_i$ appears at most once in a plan; the chain length is bounded
by $n\!\le\!20$; no branching or parallelism. User text input is
already represented as a global \texttt{[creative\_brief]} artifact, so
plans start at the first content-producing agent rather than from a
text-intake step.

\paragraph{Evaluation set.}
$209$ hand-authored goals span $43$ canonical chain shapes
(e.g.\ \emph{cinematic film}: Story$\to$Screenplay$\to$KeyFrame$\to$Video$\to$Compositor;
\emph{illustrated storytelling}: Narration$\to$Illustration$\to$Narrator$\to$Compositor;
\emph{video edit + subtitle}: IntakeVideo$\to$Transcription$\to$Translation$\to$Compositor),
distributed across nine task categories — \emph{cr} (creative
mini-drama, $38$), \emph{intake} (image-referenced, $8$), \emph{sub}
(monolingual subtitle, $15$), \emph{bilingual} (bilingual subtitle,
$8$), \emph{storytelling} (illustrated audiobook, $34$), \emph{style}
(style transfer, $8$), \emph{extend} (video extend, $8$),
\emph{highlight} (highlight reel, $6$), and \emph{complex}
(multi-modifier compositions, $84$) — and modifier flags (bilingual,
music, ambience, image-reference, longstory). The set is held-out and
is enforced to have zero overlap (under text normalization) with the
training corpus.

\paragraph{Training corpus.}
We release two training splits derived from the \emph{same}
$43$-shape blueprint as the eval set, both hand-authored under a
per-shape budget of $\max(20,\,10\!\times\!\text{eval\_count})$ so
long-tail shapes (single-eval cases) still receive $20$ training
samples. The \emph{SFT split} contains $1{,}417$ messages-format
samples with the same nine-category mix as the eval set
(\emph{complex} $539$, \emph{storytelling} $286$, \emph{cr} $253$,
\emph{sub} $79$, \emph{extend} $70$, \emph{style} $51$, \emph{intake}
$50$, \emph{bilingual} $49$, \emph{highlight} $40$); $\approx 120$
of those are \emph{long-story} samples whose user goal is a full
$\approx 2{,}500$-word original narrative, mirroring the eval's
long-input distribution. The \emph{GRPO split} contains $1{,}396$
prompts (a hand-written GRPO seed corpus plus a stratified-random
reuse of SFT user goals) with category counts \emph{complex} $519$,
\emph{storytelling} $287$, \emph{cr} $250$, \emph{sub} $80$,
\emph{extend} $71$, \emph{style} $51$, \emph{bilingual} $50$,
\emph{intake} $48$, \emph{highlight} $40$. User goals are unique
within each split and disjoint from the $209$ eval goals; details of
the SFT and GRPO learning recipes are in \S\ref{sec:sft} and
\S\ref{sec:grpo}.

\paragraph{Metrics.}
A gold chain $\mathbf{c}^*$ is a sequence in which positions may be
\emph{set-valued}: $c^*_i$ is either a single agent or a set of
interchangeable agents (e.g.\ $\{$Music, Ambience$\}$ when the audio
slot tolerates either ordering). For a predicted chain $\hat{\mathbf{c}}$:
\begin{align}
\text{chain\_acc}(\hat{\mathbf{c}},\mathbf{c}^*) &= \mathbb{1}\!\big[\hat{\mathbf{c}}\!\equiv_\text{set}\!\mathbf{c}^*\big], \\
\text{step\_acc}(\hat{\mathbf{c}},\mathbf{c}^*) &= \frac{|\{i\,:\,\hat{c}_i\in c^*_i\}|}{\max(|\hat{\mathbf{c}}|,\,|\mathbf{c}^*|)}.
\end{align}
$\equiv_\text{set}$ requires identical length and per-position
agent-set membership; the $\max$ denominator in step\_acc penalizes
length mismatches symmetrically.

\paragraph{Reference.}
Gemini 2.5 Flash on the canonical bare prompt (no fewshots, no policy
scaffolding, no topology hints) achieves chain\_acc $X.X\%$ and
step\_acc $Y.Y\%$ over the $209$ goals.\footnote{The \emph{bare} prompt
is byte-identical to the one used during training and inference of our
LoRA. Adding worked-example fewshots and topology hints lifts Gemini
to $\approx\!88\%$ chain\_acc but obscures whether the headline number
reflects model capability or template memorization, so we adopt the
bare prompt as the reference configuration.}
